% from main.tex
%%%%%%%%%%  scan idx 258  |  folio 245  %%%%%%%%%%
dérivée \(D(B)\) des \(B\)-modules, à \(M_1\otimes_{B_1}^{\mathbf L}B\otimes_{B_2}^{\mathbf L}M_2\). Les relations \(\toramp_{f_i}(E_i)\subset[a_i,b_i]\) \((i=1,2)\) impliquent \(\toramp_{B_i}(M_i)\subset[a_i,b_i]\), donc (I 5.6) \(\toramp_B(M_1\otimes_{B_1}^{\mathbf L}B\otimes_{B_2}^{\mathbf L}M_2)\subset[a,b]\), d'où l'assertion b). D'autre part, si l'on a \(\toramp_{f_1}(E_1)\subset[a_1,b_1]\), \(M_1\) est isomorphe à un complexe de \(B_1\)-modules plats, nuls en dehors de \([a_1,b_1]\), et par suite si \(E_2\) (donc \(M_2\)) est acyclique en dehors de \([a_2,b_2]\), \(M_1\otimes_{B_1}^{\mathbf L}B\otimes_{B_2}^{\mathbf L}M_2\) est acyclique en dehors de \([a,b]\), d'où l'assertion a), ce qui achève la démonstration.

\textbf{Corollaire 3.5.1.} Sous les hypothèses de (3.5), si \(f_1\) et \(f_2\) sont de tor-dimension finie (3.1), il en est de même de \(f\).

\textbf{Corollaire 3.5.2.} Soit
\[
\begin{tikzcd}
X \arrow[d,"f"'] & X' \arrow[l,"g"'] \arrow[d,"f'"]\\
Y & Y' \arrow[l,"h"']
\end{tikzcd}
\]
un carré cartésien de schémas, où \(X\) et \(Y'\) sont supposés tor-indépendants sur \(Y\), et soit \(E\in\Ob D(X)\). Alors on a
\[
  \toramp_f(E)\subset[a,b]\Longrightarrow
  \toramp_{f'}(\mathbf Lg^*(E))\subset[a,b],
\]
et par suite le foncteur \(\mathbf Lg^*:D^{-}(X)\to D^{-}(X')\) induit un foncteur
\[
  \mathbf Lg^*:D(f)_{\mathrm{torf}}\longrightarrow D(f')_{\mathrm{torf}} .
\]



% Additions for Exposé III, pages 253--280
\renewcommand{\Ob}{\operatorname{Ob}}
\providecommand{\D}{\mathrm D}
\providecommand{\Ho}{\mathrm H}
\providecommand{\RHom}{\mathrm R\!\operatorname{Hom}}
\providecommand{\parfamp}{\operatorname{parf.amp}}
\providecommand{\tamp}{\operatorname{tor.amp}}
\providecommand{\Qcoh}{\operatorname{Qcoh}}
\providecommand{\Mod}{\operatorname{Mod}}
\providecommand{\Coh}{\operatorname{Coh}}
\providecommand{\Lib}{\operatorname{Lib}}
\providecommand{\Loclib}{\operatorname{Loclib}}
\providecommand{\qcoh}{\mathrm{qcoh}}
\providecommand{\coh}{\mathrm{coh}}
\providecommand{\parf}{\mathrm{parf}}
\providecommand{\torf}{\mathrm{torf}}
\renewcommand{\id}{\operatorname{id}}
\providecommand{\uparf}{\underline{\mathrm{parf}}\rule[-1ex]{0pt}{0pt}}
\providecommand{\Coker}{\operatorname{Coker}}
\providecommand{\rank}{\operatorname{rank}}
\providecommand{\rg}{\operatorname{rg}}
\providecommand{\Hom}{\operatorname{Hom}}
\providecommand{\Ext}{\operatorname{Ext}}
\providecommand{\Spec}{\operatorname{Spec}}
\providecommand{\Cone}{\operatorname{Cone}}
\providecommand{\pt}{\operatorname{Pt}}
\providecommand{\ind}{\operatorname{ind}}
\hbadness=10000
\hfuzz=20pt
\emergencystretch=3em

